% ===========================================================================
%  evnx for DevOps / Platform Engineers  —  K8s, Terraform, GitHub Actions
%  Persona NK.  Build: pdflatex -interaction=nonstopmode devops.tex
% ===========================================================================
\documentclass[aspectratio=169,10pt]{beamer}

\newcommand{\capturedir}{../captures}
% ---------------------------------------------------------------------------
%  Shared preamble for the seven audience decks.
%
%  Each deck does, before the document body starts:
%     \newcommand{\capturedir}{../captures}
%     % ---------------------------------------------------------------------------
%  Shared preamble for the seven audience decks.
%
%  Each deck does, before the document body starts:
%     \newcommand{\capturedir}{../captures}
%     % ---------------------------------------------------------------------------
%  Shared preamble for the seven audience decks.
%
%  Each deck does, before the document body starts:
%     \newcommand{\capturedir}{../captures}
%     \input{persona-preamble}
%
%  \capturedir points at ../captures so both pools resolve:
%     \capture{24-scan.txt}              — the original v0.5.0 captures
%     \capture{persona/fe-01-scan-dist.txt} — captures recorded for these decks
% ---------------------------------------------------------------------------

\input{../common/evnx-preamble}

% ---- provenance ------------------------------------------------------------
% The original decks were pinned to tag v0.5.0 (85a52c3). Everything captured
% for the persona decks was re-run against 0.5.2, so the footline differs.
\renewcommand{\verifiedon}{\tiny\textcolor{evnxMuted}{verified against evnx 0.5.2 --- real run, not transcribed}}

% ---- audience badge on the title slide ------------------------------------
\newcommand{\audience}[1]{%
  \begin{center}\small\textcolor{evnxAccent}{\textbf{Audience: #1}}\end{center}}

% ---- verdict chips ---------------------------------------------------------
\newcommand{\solved}{\textcolor{evnxOk}{\ensuremath{\checkmark}}}
\newcommand{\partialmark}{\textcolor{evnxWarn}{\textbf{\textasciitilde}}}
\newcommand{\unsolved}{\textcolor{evnxErr}{\ensuremath{\times}}}

% ---- a correction block: something the source material got wrong -----------
\newenvironment{correction}[1]{%
  \begin{alertblock}{\small Correction --- #1}\small}{\end{alertblock}}

% ---- a trap block: works in the demo, fails in real use --------------------
\newenvironment{trap}[1]{%
  \setbeamercolor{block title}{bg=evnxWarn,fg=white}%
  \setbeamercolor{block body}{bg=evnxWarn!8}%
  \begin{block}{\small #1}\small}{\end{block}}

% ---- the persona card ------------------------------------------------------
\newcommand{\personacard}[4]{%
  \begin{block}{#1}
    \small
    \textbf{Stack:} #2 \\
    \textbf{Ships:} #3 \\
    \textbf{Loses time to:} #4
  \end{block}}

% ---- speaker note (rendered small and muted, not hidden) -------------------
\newcommand{\saythis}[1]{%
  \vspace{0.4em}{\footnotesize\textcolor{evnxMuted}{\textit{Say:} #1}}}

%
%  \capturedir points at ../captures so both pools resolve:
%     \capture{24-scan.txt}              — the original v0.5.0 captures
%     \capture{persona/fe-01-scan-dist.txt} — captures recorded for these decks
% ---------------------------------------------------------------------------

% ---------------------------------------------------------------------------
%  evnx presentation preamble — shared by all three decks
%
%  Engine: pdfLaTeX (also builds under XeLaTeX / LuaLaTeX).
%  Packages: beamer, listings, xcolor, amssymb, newunicodechar, booktabs,
%            fancyvrb, hyperref — all in texlive-latex-recommended + beamer.
%
%  \capturedir must be defined by the including document BEFORE \input of
%  this file, e.g.  \newcommand{\capturedir}{../captures}
% ---------------------------------------------------------------------------

\usetheme{Madrid}
\usecolortheme{seahorse}
\usefonttheme{professionalfonts}
\setbeamertemplate{navigation symbols}{}
\setbeamertemplate{caption}[numbered]
\setbeamerfont{frametitle}{size=\large}
\setbeamerfont{title}{size=\LARGE,series=\bfseries}

\usepackage[T1]{fontenc}
\usepackage[utf8]{inputenc}
\usepackage{lmodern}
\usepackage{amssymb}
\usepackage{booktabs}
\usepackage{xcolor}
\usepackage{listings}
\usepackage{fancyvrb}
\usepackage{newunicodechar}
\usepackage{tikz}
\usetikzlibrary{arrows.meta,positioning,shapes.geometric,fit,backgrounds}

% ---- palette ---------------------------------------------------------------
\definecolor{evnxInk}{HTML}{1F2933}
\definecolor{evnxAccent}{HTML}{2F6F4E}
\definecolor{evnxWarn}{HTML}{B7791F}
\definecolor{evnxErr}{HTML}{B03A2E}
\definecolor{evnxOk}{HTML}{2F6F4E}
\definecolor{evnxMuted}{HTML}{6B7280}
\definecolor{evnxTermBg}{HTML}{F5F5F2}
\definecolor{evnxRule}{HTML}{D1D5DB}

\setbeamercolor{structure}{fg=evnxAccent}
\setbeamercolor{alerted text}{fg=evnxErr}

% ---- glyph mapping for body text ------------------------------------------
% evnx emits exactly these 11 non-ASCII characters (enumerated from the
% captured logs, not guessed).
\newcommand{\evnxCheck}{\textcolor{evnxOk}{\ensuremath{\checkmark}}}
\newcommand{\evnxCross}{\textcolor{evnxErr}{\ensuremath{\times}}}
\newcommand{\evnxWarnSign}{\textcolor{evnxWarn}{\textbf{!}}}

\newunicodechar{✓}{\evnxCheck}
\newunicodechar{✗}{\evnxCross}
\newunicodechar{⚠}{\evnxWarnSign}
\newunicodechar{→}{\ensuremath{\rightarrow}}
\newunicodechar{↔}{\ensuremath{\leftrightarrow}}
\newunicodechar{─}{\textendash}
\newunicodechar{—}{---}
\newunicodechar{·}{\textperiodcentered}
\newunicodechar{…}{\ldots}
\newunicodechar{•}{\textbullet}
\newunicodechar{️}{}                % U+FE0F variation selector — invisible

% ---- terminal listing style ------------------------------------------------
% `literate` is required: newunicodechar does not apply inside listings.
\lstdefinestyle{evnxterm}{
  basicstyle=\ttfamily\scriptsize,
  backgroundcolor=\color{evnxTermBg},
  frame=single,
  rulecolor=\color{evnxRule},
  framesep=4pt,
  breaklines=true,
  breakatwhitespace=false,
  columns=fullflexible,
  keepspaces=true,
  showstringspaces=false,
  upquote=true,
  extendedchars=true,
  literate=
    {✓}{{\textcolor{evnxOk}{$\checkmark$}}}1
    {✗}{{\textcolor{evnxErr}{$\times$}}}1
    {⚠}{{\textcolor{evnxWarn}{\textbf{!}}}}1
    {→}{{$\rightarrow$}}1
    {↔}{{$\leftrightarrow$}}1
    {─}{{-}}1
    {—}{{---}}1
    {·}{{$\cdot$}}1
    {…}{{\ldots}}1
    {•}{{\textbullet}}1
    {️}{{}}0
    {é}{{\'e}}1
}

\lstdefinestyle{evnxcode}{
  style=evnxterm,
  basicstyle=\ttfamily\scriptsize,
  backgroundcolor=\color{white},
}

\lstset{style=evnxterm}

% ---- helper macros ---------------------------------------------------------
% \capture{file}        — include a real captured run verbatim
% \capturepart{f}{a}{b} — include lines a..b of a captured run
\newcommand{\capture}[1]{\lstinputlisting[style=evnxterm]{\capturedir/#1}}
\newcommand{\capturepart}[3]{%
  \lstinputlisting[style=evnxterm,firstline=#2,lastline=#3]{\capturedir/#1}}

% a labelled "real output" box
\newcommand{\realrun}[1]{%
  \begin{center}\scriptsize\textcolor{evnxMuted}{real run \textendash\ #1}\end{center}}

\newcommand{\cmd}[1]{\texttt{\textbf{#1}}}
\newcommand{\flag}[1]{\texttt{#1}}
\newcommand{\file}[1]{\texttt{#1}}

% exit-code chips
\newcommand{\exitzero}{\textcolor{evnxOk}{\texttt{0}}}
\newcommand{\exitone}{\textcolor{evnxWarn}{\texttt{1}}}
\newcommand{\exittwo}{\textcolor{evnxErr}{\texttt{2}}}

% a standard "verified against" footline for factual slides
\newcommand{\verifiedon}{\tiny\textcolor{evnxMuted}{verified against evnx 0.5.0 @ 9cfe75b}}

\usepackage[colorlinks=true,linkcolor=evnxAccent,urlcolor=evnxAccent]{hyperref}


% ---- provenance ------------------------------------------------------------
% The original decks were pinned to tag v0.5.0 (85a52c3). Everything captured
% for the persona decks was re-run against 0.5.2, so the footline differs.
\renewcommand{\verifiedon}{\tiny\textcolor{evnxMuted}{verified against evnx 0.5.2 --- real run, not transcribed}}

% ---- audience badge on the title slide ------------------------------------
\newcommand{\audience}[1]{%
  \begin{center}\small\textcolor{evnxAccent}{\textbf{Audience: #1}}\end{center}}

% ---- verdict chips ---------------------------------------------------------
\newcommand{\solved}{\textcolor{evnxOk}{\ensuremath{\checkmark}}}
\newcommand{\partialmark}{\textcolor{evnxWarn}{\textbf{\textasciitilde}}}
\newcommand{\unsolved}{\textcolor{evnxErr}{\ensuremath{\times}}}

% ---- a correction block: something the source material got wrong -----------
\newenvironment{correction}[1]{%
  \begin{alertblock}{\small Correction --- #1}\small}{\end{alertblock}}

% ---- a trap block: works in the demo, fails in real use --------------------
\newenvironment{trap}[1]{%
  \setbeamercolor{block title}{bg=evnxWarn,fg=white}%
  \setbeamercolor{block body}{bg=evnxWarn!8}%
  \begin{block}{\small #1}\small}{\end{block}}

% ---- the persona card ------------------------------------------------------
\newcommand{\personacard}[4]{%
  \begin{block}{#1}
    \small
    \textbf{Stack:} #2 \\
    \textbf{Ships:} #3 \\
    \textbf{Loses time to:} #4
  \end{block}}

% ---- speaker note (rendered small and muted, not hidden) -------------------
\newcommand{\saythis}[1]{%
  \vspace{0.4em}{\footnotesize\textcolor{evnxMuted}{\textit{Say:} #1}}}

%
%  \capturedir points at ../captures so both pools resolve:
%     \capture{24-scan.txt}              — the original v0.5.0 captures
%     \capture{persona/fe-01-scan-dist.txt} — captures recorded for these decks
% ---------------------------------------------------------------------------

% ---------------------------------------------------------------------------
%  evnx presentation preamble — shared by all three decks
%
%  Engine: pdfLaTeX (also builds under XeLaTeX / LuaLaTeX).
%  Packages: beamer, listings, xcolor, amssymb, newunicodechar, booktabs,
%            fancyvrb, hyperref — all in texlive-latex-recommended + beamer.
%
%  \capturedir must be defined by the including document BEFORE \input of
%  this file, e.g.  \newcommand{\capturedir}{../captures}
% ---------------------------------------------------------------------------

\usetheme{Madrid}
\usecolortheme{seahorse}
\usefonttheme{professionalfonts}
\setbeamertemplate{navigation symbols}{}
\setbeamertemplate{caption}[numbered]
\setbeamerfont{frametitle}{size=\large}
\setbeamerfont{title}{size=\LARGE,series=\bfseries}

\usepackage[T1]{fontenc}
\usepackage[utf8]{inputenc}
\usepackage{lmodern}
\usepackage{amssymb}
\usepackage{booktabs}
\usepackage{xcolor}
\usepackage{listings}
\usepackage{fancyvrb}
\usepackage{newunicodechar}
\usepackage{tikz}
\usetikzlibrary{arrows.meta,positioning,shapes.geometric,fit,backgrounds}

% ---- palette ---------------------------------------------------------------
\definecolor{evnxInk}{HTML}{1F2933}
\definecolor{evnxAccent}{HTML}{2F6F4E}
\definecolor{evnxWarn}{HTML}{B7791F}
\definecolor{evnxErr}{HTML}{B03A2E}
\definecolor{evnxOk}{HTML}{2F6F4E}
\definecolor{evnxMuted}{HTML}{6B7280}
\definecolor{evnxTermBg}{HTML}{F5F5F2}
\definecolor{evnxRule}{HTML}{D1D5DB}

\setbeamercolor{structure}{fg=evnxAccent}
\setbeamercolor{alerted text}{fg=evnxErr}

% ---- glyph mapping for body text ------------------------------------------
% evnx emits exactly these 11 non-ASCII characters (enumerated from the
% captured logs, not guessed).
\newcommand{\evnxCheck}{\textcolor{evnxOk}{\ensuremath{\checkmark}}}
\newcommand{\evnxCross}{\textcolor{evnxErr}{\ensuremath{\times}}}
\newcommand{\evnxWarnSign}{\textcolor{evnxWarn}{\textbf{!}}}

\newunicodechar{✓}{\evnxCheck}
\newunicodechar{✗}{\evnxCross}
\newunicodechar{⚠}{\evnxWarnSign}
\newunicodechar{→}{\ensuremath{\rightarrow}}
\newunicodechar{↔}{\ensuremath{\leftrightarrow}}
\newunicodechar{─}{\textendash}
\newunicodechar{—}{---}
\newunicodechar{·}{\textperiodcentered}
\newunicodechar{…}{\ldots}
\newunicodechar{•}{\textbullet}
\newunicodechar{️}{}                % U+FE0F variation selector — invisible

% ---- terminal listing style ------------------------------------------------
% `literate` is required: newunicodechar does not apply inside listings.
\lstdefinestyle{evnxterm}{
  basicstyle=\ttfamily\scriptsize,
  backgroundcolor=\color{evnxTermBg},
  frame=single,
  rulecolor=\color{evnxRule},
  framesep=4pt,
  breaklines=true,
  breakatwhitespace=false,
  columns=fullflexible,
  keepspaces=true,
  showstringspaces=false,
  upquote=true,
  extendedchars=true,
  literate=
    {✓}{{\textcolor{evnxOk}{$\checkmark$}}}1
    {✗}{{\textcolor{evnxErr}{$\times$}}}1
    {⚠}{{\textcolor{evnxWarn}{\textbf{!}}}}1
    {→}{{$\rightarrow$}}1
    {↔}{{$\leftrightarrow$}}1
    {─}{{-}}1
    {—}{{---}}1
    {·}{{$\cdot$}}1
    {…}{{\ldots}}1
    {•}{{\textbullet}}1
    {️}{{}}0
    {é}{{\'e}}1
}

\lstdefinestyle{evnxcode}{
  style=evnxterm,
  basicstyle=\ttfamily\scriptsize,
  backgroundcolor=\color{white},
}

\lstset{style=evnxterm}

% ---- helper macros ---------------------------------------------------------
% \capture{file}        — include a real captured run verbatim
% \capturepart{f}{a}{b} — include lines a..b of a captured run
\newcommand{\capture}[1]{\lstinputlisting[style=evnxterm]{\capturedir/#1}}
\newcommand{\capturepart}[3]{%
  \lstinputlisting[style=evnxterm,firstline=#2,lastline=#3]{\capturedir/#1}}

% a labelled "real output" box
\newcommand{\realrun}[1]{%
  \begin{center}\scriptsize\textcolor{evnxMuted}{real run \textendash\ #1}\end{center}}

\newcommand{\cmd}[1]{\texttt{\textbf{#1}}}
\newcommand{\flag}[1]{\texttt{#1}}
\newcommand{\file}[1]{\texttt{#1}}

% exit-code chips
\newcommand{\exitzero}{\textcolor{evnxOk}{\texttt{0}}}
\newcommand{\exitone}{\textcolor{evnxWarn}{\texttt{1}}}
\newcommand{\exittwo}{\textcolor{evnxErr}{\texttt{2}}}

% a standard "verified against" footline for factual slides
\newcommand{\verifiedon}{\tiny\textcolor{evnxMuted}{verified against evnx 0.5.0 @ 9cfe75b}}

\usepackage[colorlinks=true,linkcolor=evnxAccent,urlcolor=evnxAccent]{hyperref}


% ---- provenance ------------------------------------------------------------
% The original decks were pinned to tag v0.5.0 (85a52c3). Everything captured
% for the persona decks was re-run against 0.5.2, so the footline differs.
\renewcommand{\verifiedon}{\tiny\textcolor{evnxMuted}{verified against evnx 0.5.2 --- real run, not transcribed}}

% ---- audience badge on the title slide ------------------------------------
\newcommand{\audience}[1]{%
  \begin{center}\small\textcolor{evnxAccent}{\textbf{Audience: #1}}\end{center}}

% ---- verdict chips ---------------------------------------------------------
\newcommand{\solved}{\textcolor{evnxOk}{\ensuremath{\checkmark}}}
\newcommand{\partialmark}{\textcolor{evnxWarn}{\textbf{\textasciitilde}}}
\newcommand{\unsolved}{\textcolor{evnxErr}{\ensuremath{\times}}}

% ---- a correction block: something the source material got wrong -----------
\newenvironment{correction}[1]{%
  \begin{alertblock}{\small Correction --- #1}\small}{\end{alertblock}}

% ---- a trap block: works in the demo, fails in real use --------------------
\newenvironment{trap}[1]{%
  \setbeamercolor{block title}{bg=evnxWarn,fg=white}%
  \setbeamercolor{block body}{bg=evnxWarn!8}%
  \begin{block}{\small #1}\small}{\end{block}}

% ---- the persona card ------------------------------------------------------
\newcommand{\personacard}[4]{%
  \begin{block}{#1}
    \small
    \textbf{Stack:} #2 \\
    \textbf{Ships:} #3 \\
    \textbf{Loses time to:} #4
  \end{block}}

% ---- speaker note (rendered small and muted, not hidden) -------------------
\newcommand{\saythis}[1]{%
  \vspace{0.4em}{\footnotesize\textcolor{evnxMuted}{\textit{Say:} #1}}}


\title{evnx for platform engineers}
\subtitle{A CI hygiene tool that is honest about not being a secret store}
\date{}

\begin{document}

\begin{frame}[plain]
  \titlepage
  \audience{GitHub Actions · Kubernetes · Helm · Argo · Terraform · AWS SM + ESO}
  \begin{center}\footnotesize
  \textcolor{evnxMuted}{evnx 0.5.2. Every terminal block is a real run.}
  \end{center}
\end{frame}

% =========================================================================
\begin{frame}{This deck in one slide}
  \begin{block}{The recommendation, up front}
    Adopt the \textbf{CLI} across all 40 repos as a CI gate. Use \textbf{evnx
    cloud} for developer-local environments only. Do \textbf{not} plan to
    replace AWS Secrets Manager with it this year.
  \end{block}

  \vspace{0.7em}
  \begin{block}{Why that split}
    The things you need from a hygiene tool --- a container image, a real exit
    contract, stable SARIF, many output formats --- are all finished. The
    things you need from a \emph{platform} secret store --- machine
    identities, org policy, SSO, rotation, ESO integration --- are not
    started. Pretending otherwise wastes your quarter.
  \end{block}

  \saythis{Leading with the limits is what gets this audience to listen to the rest.}
\end{frame}

\begin{frame}{Contents}\tableofcontents[hideallsubsections]\end{frame}

% =========================================================================
\section{Before lunch}

\begin{frame}{Three tickets and an image}
  \small
  \begin{block}{The tickets}
    ``Add \texttt{FEATURE\_X\_URL} to staging.'' ``Prod is missing
    \texttt{SENTRY\_DSN}.'' ``Rotate the Stripe key everywhere.''
  \end{block}

  \begin{block}{The rotation, in full}
    Update AWS SM → re-sync ExternalSecrets → restart pods → update two GitHub
    environments → discover the key is \emph{also} hard-coded in a Helm values
    file from 2023.
  \end{block}

  \begin{block}{The afternoon}
    An image in GHCR turns out to contain \file{.env}. The service's Dockerfile
    does \texttt{COPY . .} with no \file{.dockerignore}.
  \end{block}

  \vspace{0.4em}
  {\footnotesize evnx addresses the \textbf{last} one properly, the first two
  partially, and the rotation not at all. The rest of this deck is which is which.}
\end{frame}

% =========================================================================
\section{The CI primitives}

\begin{frame}{No install step}
  \begin{center}\large\ttfamily
  docker run -{}-rm -v "\$PWD:/work" ghcr.io/urwithajit9/evnx scan
  \end{center}

  \vspace{1em}
  \begin{itemize}
    \item One image, pinnable by tag, usable as a \texttt{container:} in a
          reusable workflow --- no \texttt{setup-} action, no toolchain, no
          cache key.
    \item The prebuilt binaries are built with \texttt{-{}-all-features}, so
          \texttt{migrate}, \texttt{backup} and \texttt{cloud} are all present.
          Only \cmd{cargo install evnx} compiles with \texttt{default = []}
          and gives you 11 of 17 commands.
    \item Also on Homebrew, Scoop, winget, npm, PyPI and as a Marketplace
          action.
  \end{itemize}

  \vspace{0.7em}
  {\footnotesize Check the GHCR package is public before you templatise it
  across the org.}
\end{frame}

\begin{frame}[fragile]{The exit contract --- the slide that matters most}
  \begin{center}
  \begin{tabular}{@{}cll@{}}
    \toprule
    \textbf{Code} & \textbf{Means} & \textbf{Your gate should} \\
    \midrule
    \exitzero & clean                  & continue \\
    \exitone  & a finding              & fail, and surface it \\
    \exittwo  & \textbf{could not run} & fail --- \emph{this is not a pass} \\
    \bottomrule
  \end{tabular}
  \end{center}

  \vspace{0.6em}
  \capturepart{persona/ops-06-exit2.txt}{1}{9}
  \realrun{a path that does not exist}

  \vspace{0.3em}
  {\footnotesize \textbf{Action item:} audit any pipeline that treats every
  non-zero as ``finding''. Under the old behaviour, \exittwo\ conditions
  sometimes returned \exitzero\ --- a gate that could not run looked green.}
\end{frame}

\begin{frame}[fragile]{Annotations on the right lines}
  \capture{persona/ops-01-validate-github.txt}
  \realrun{\texttt{evnx validate -{}-format github}}

  \vspace{0.3em}
  {\footnotesize Workflow-command format, correct \texttt{file=} and
  \texttt{line=}, \texttt{\%0A} for multi-line suggestions. Annotations land on
  the diff, so a reviewer sees the finding without opening the log.}
\end{frame}

\begin{frame}[fragile]{SARIF that does not reset your alerts}
  \capturepart{persona/ops-02-scan-sarif.txt}{1}{31}
  \realrun{\texttt{evnx scan . -{}-format sarif}}

  \vspace{0.3em}
  {\footnotesize Stable \texttt{ruleId}, a \texttt{rules[]} block with
  \texttt{helpUri}, and \texttt{partialFingerprints} keyed on
  \texttt{file::variable::rule}. That fingerprint is what stops GitHub code
  scanning from closing and reopening every alert on each run. Expect
  \textbf{one} alert reset when you upgrade to it, then stability.}
\end{frame}

\begin{frame}[fragile]{Stream separation --- why the pipe above is safe}
  \begin{block}{The rule}
    Machine formats (\texttt{json}, \texttt{sarif}, \texttt{github}) go to
    \textbf{stdout}. Human trailers --- the ``Docs:'' line, next steps,
    summaries --- go to \textbf{stderr}.
  \end{block}

  \vspace{0.7em}
\begin{lstlisting}[style=evnxcode]
evnx scan . --format sarif > evnx.sarif || test $? -eq 1
\end{lstlisting}

  \vspace{0.5em}
  {\footnotesize The redirect captures clean SARIF. The \texttt{|| test \$? -eq 1}
  lets findings upload as alerts while still failing the job on \exittwo\ ---
  adjust to your policy, but make the \exittwo\ case explicit.}
\end{frame}

\begin{frame}[fragile]{The reusable workflow}
\begin{lstlisting}[style=evnxcode]
on: [pull_request]
jobs:
  env-hygiene:
    runs-on: ubuntu-latest
    container: ghcr.io/urwithajit9/evnx:0.5.0
    steps:
      - uses: actions/checkout@v4
      - run: evnx doctor --format json
      - run: evnx sync --check
      - run: evnx validate --format github
      - run: evnx scan . --format sarif > evnx.sarif || test $? -eq 1
      - uses: github/codeql-action/upload-sarif@v3
        if: always()
        with: { sarif_file: evnx.sarif }
\end{lstlisting}
  {\footnotesize Drop this in an org template repo. It is the whole A-stage
  answer.}
\end{frame}

% =========================================================================
\section{The failure mode to know about}

\begin{frame}[fragile]{A scanner's worst outcome is a confident pass}
  \capture{persona/fe-01-scan-dist.txt}
  \realrun{\file{dist/bundle.js} contains a live \texttt{sk\_live\_} key}

  \vspace{0.4em}
  \begin{trap}{\file{dist/} and \file{build/} are on the default exclusion list}
    Excluding build output is correct when scanning a repo. It is wrong when
    the excluded directory is the path you \emph{asked} for --- and evnx says
    ``✓ No secrets detected'', ``0 files scanned'', exit \exitzero.
    \file{src/commands/scan/filters.rs:106}. Tracked \textbf{NEW-1, P0}.
  \end{trap}

  \vspace{0.3em}
  {\footnotesize Same for \file{.ipynb}, which is not in the scannable
  extension allowlist at all (\textbf{NEW-2}).}
\end{frame}

\begin{frame}[fragile]{Assert the denominator, not just the exit code}
  \capture{persona/ops-09-files-scanned.txt}
  \vspace{-0.2em}
\begin{lstlisting}[style=evnxcode]
# fail if the scan looked at nothing
n=$(evnx scan "$TARGET" --format json | jq '.files_scanned')
test "$n" -gt 0 || { echo "scan covered no files"; exit 1; }
\end{lstlisting}

  \vspace{0.8em}
  \begin{block}{Generalise it}
    This is not an evnx-specific hack. Any coverage-based gate --- linters,
    scanners, test runners --- can pass by looking at nothing. Assert the
    denominator in every one of them.
  \end{block}

  \saythis{This slide travels well beyond evnx, which is why it is worth
  keeping even after NEW-1 is fixed.}
\end{frame}

% =========================================================================
\section{Conversion and migration}

\begin{frame}[fragile]{\cmd{convert} --- 14 targets, all plaintext}
  \capture{persona/ops-03-convert-tf.txt}
  \realrun{\texttt{evnx convert -{}-to terraform}}

  \vspace{0.3em}
  {\footnotesize Also: \texttt{json, yaml, shell, aws-secrets, gcp-secrets,
  azure-keyvault, github-actions, docker-compose, kubernetes, doppler, heroku,
  vercel, railway}. \textbf{Pipe, never commit.}}

  \vspace{0.4em}
  \begin{trap}{The Terraform caveat you already know}
    \texttt{TF\_VAR\_*} and tfvars mean the value ends up in
    \textbf{Terraform state}, in plaintext. evnx converting the format does not
    change that. There is no Terraform provider for evnx.
  \end{trap}
\end{frame}

\begin{frame}[fragile]{\cmd{migrate} is a command generator}
  \capturepart{persona/ops-08-migrate-live-aws.txt}{1}{19}
  \realrun{\texttt{evnx migrate -{}-to aws-secrets-manager} --- \textbf{no} \flag{-{}-dry-run}}

  \vspace{0.3em}
  \begin{correction}{eight of nine destinations contact nothing}
    They print the platform's own CLI commands so an operator can review and
    audit before running them. Only \texttt{github-actions} actually transfers
    (sealed box → REST). Until v0.5.0 the summary said ``\texttt{N uploaded}''
    for all nine, which was misleading; the banner now reads ``Generate
    destination commands''.
  \end{correction}
\end{frame}

\begin{frame}[fragile]{Two \cmd{migrate} behaviours worth relying on}
  \capturepart{persona/ops-04-migrate-dry.txt}{1}{12}
  \realrun{\flag{-{}-dry-run} \emph{without} \flag{-{}-secret-name}}

  \vspace{0.4em}
  \begin{itemize}\footnotesize
    \item \flag{-{}-dry-run} still demands the required arguments --- rather
          than printing a command with a placeholder where the name goes.
    \item A non-interactive run must pass \flag{-{}-to}. It no longer guesses
          the destination.
    \item \flag{-{}-skip-existing} / \flag{-{}-overwrite} are \textbf{inert}
          for the eight command-generating destinations. No error if you pass
          them.
  \end{itemize}
\end{frame}

% =========================================================================
\section{Where it is not a platform tool}

\begin{frame}{The five gaps, stated plainly}
  \small
  \begin{tabular}{@{}clp{6.4cm}@{}}
    \toprule
    & \textbf{Gap} & \textbf{Consequence for you} \\
    \midrule
    \unsolved & \textbf{Machine identities} & CI pulling from a vault needs the
      \emph{master password} plus a token. That is a human credential in a
      pipeline. This is the single blocker for broad cloud adoption \\
    \unsolved & Rotation \& dynamic secrets & no leases, no short-lived DB
      credentials. Keep AWS SM rotation \\
    \unsolved & Terraform provider          & no IaC-managed vaults \\
    \partialmark & Org policy \& audit export & per-vault audit log and
      \cmd{cloud history} exist; no SIEM export, SSO/SCIM or admin console \\
    \partialmark & Offboarding across vaults & \cmd{vault revoke} re-keys the
      vault properly, but there is no ``remove user everywhere'' \\
    \bottomrule
  \end{tabular}

  \vspace{0.6em}
  \saythis{``If someone tells you this replaces Vault or Secrets Manager this
  year, they have not read the machine-identity row.''}
\end{frame}

\begin{frame}{And two you might expect but should not}
  \small
  \begin{block}{Image layers}
    \cmd{evnx scan .} runs \emph{before} \texttt{docker build}. It does not
    inspect image layers and \cmd{doctor} does not check
    \file{.dockerignore} --- so the GHCR incident from slide 5 is caught only
    if the \file{.env} is also in the working tree at scan time.
  \end{block}

  \begin{block}{Git history}
    \cmd{scan} looks at the working tree. Historical commits are not its focus
    --- pair it with gitleaks or trufflehog for history, and use evnx as the
    pre-commit and PR gate that stops new ones.
  \end{block}

  \vspace{0.4em}
  {\footnotesize Both are honest scoping decisions rather than oversights, but
  neither is obvious from the README.}
\end{frame}

\begin{frame}{Two sources of truth, and how to keep it sane}
  \begin{block}{The split that works}
    \begin{tabular}{@{}ll@{}}
      \textbf{Developer-local} & evnx vaults are the source of truth \\
      \textbf{Staging / prod}  & AWS SM remains the source of truth \\
      \textbf{The contract between them} & \file{.env.example}, checked in CI \\
    \end{tabular}
  \end{block}

  \vspace{0.8em}
  \begin{block}{Closing the ``new variable is missing in prod'' gap today}
    There is no direct ``compare vault vs AWS SM''. Export the target to a
    file and \cmd{evnx diff} against it --- or enforce \file{.env.example} as
    the contract with \cmd{sync -{}-check} on every PR, which catches the
    variable before it reaches a deploy.
  \end{block}
\end{frame}

% =========================================================================
\section{Adoption}

\begin{frame}[fragile]{Rollout in three steps}
\begin{lstlisting}[style=evnxcode]
# 1. org template repo -- the reusable workflow, pinned by tag
#    (slide 10). Roll to 40 repos via a required check.

# 2. pre-commit, same three hooks everywhere
repos:
  - repo: https://github.com/urwithajit9/evnx
    rev: v0.5.0
    hooks: [{id: evnx-scan}, {id: evnx-validate}, {id: evnx-diff}]

# 3. shared detector rules, committed
#    .evnx.toml
[[scan.patterns]]
name  = "Acme internal token"
regex = "ACME-[A-Z0-9]{32}"
\end{lstlisting}
  {\footnotesize Step 3 is what makes the gate org-specific rather than
  generic --- and it is committed, so it reviews like code.}
\end{frame}

\begin{frame}{The scorecard, honestly}
  \begin{center}
  \begin{tabular}{@{}clp{6.2cm}@{}}
    \toprule
    & \textbf{Count} & \textbf{Items} \\
    \midrule
    \solved     & 6 & container image, exit contract, stable SARIF, all convert targets, bulk migrate, org-wide secret scanning \\
    \partialmark & 8 & per-repo standards, target drift, developer/platform split, Helm values, image contents, audit, offboarding, Terraform vars \\
    \unsolved   & 3 & \textbf{machine identities}, rotation, Terraform provider \\
    \bottomrule
  \end{tabular}
  \end{center}

  \vspace{0.7em}
  \begin{block}{One sentence}
    A solid CI hygiene and conversion tool. Not yet a platform secret store ---
    and the roadmap agrees with that assessment rather than arguing with it.
  \end{block}
\end{frame}

\begin{frame}{What is coming, in priority order}
  \small
  \begin{tabular}{@{}clp{5.6cm}@{}}
    \toprule
    & \textbf{Feature} & \textbf{Unblocks} \\
    \midrule
    P0 & Scan an explicitly-named excluded path & the confident-pass failure \\
    1  & Machine identities / service accounts (vault-scoped, no master password) & CI and cluster adoption \\
    2  & Org policy: shared config, required checks, central view & 40 repos, one standard \\
    3  & \cmd{diff -{}-against aws:…\textbar k8s:ns/secret\textbar github:env} & ``missing in prod'' \\
    4  & External Secrets Operator provider & the two-sources-of-truth split \\
    5  & \cmd{scan -{}-image ghcr.io/…} + \file{.dockerignore} in \cmd{doctor} & the GHCR incident \\
    6  & SSO/SCIM, audit export, org-wide offboarding & compliance \\
    \bottomrule
  \end{tabular}
\end{frame}

\begin{frame}[plain]
  \vfill
  \begin{center}
    {\LARGE\textbf{Questions}}\\[1.2em]
    {\small \url{https://www.evnx.dev/guides} · \url{https://github.com/urwithajit9/evnx}}\\[0.8em]
    {\footnotesize\textcolor{evnxMuted}{evnx 0.5.2 · every terminal block is a real run}}
  \end{center}
  \vfill
\end{frame}

\end{document}
